% Source-aligned English for printed p.218.
\clearpage
\SourcePage{218}

\LayerTarget{tohoku:transition:degeneration-case}
\LayerTarget{tohoku:corollary:theorem-5.6.3-locally-free}
\noindent\textsc{Corollary.}\quad \textit{If the $G$-$O$-module $A$ is
locally isomorphic, as an $O$-module, to $O^n$, then there are canonical
isomorphisms}

\LayerTarget{tohoku:equation:5.6.12}
\begin{equation}
\operatorname{Ext}_{O,G}^n(X;A,B)
  =H^n\bigl(X;G,\mathbf{Hom}_O(A,B)\bigr).
\tag{5.6.12}
\end{equation}

\LayerTarget{tohoku:proof:theorem-5.6.3-locally-free-corollary}
This follows from spectral sequence III, since then
$\mathbf{Ext}_O^q(A,B)=0$ for $q>0$ (Prop. 4.2.3). In particular, taking
$A=O$, one obtains:

\LayerTarget{tohoku:corollary:theorem-5.6.3-structure-sheaf}
\noindent\textsc{Corollary 2.}\quad \textit{For every $G$-$O$-module $B$
there are functorial isomorphisms:}

\LayerTarget{tohoku:equation:5.6.14}
\begin{equation}
\operatorname{Ext}_{O,G}^n(X;O,B)=H^n(X;G,B).
\tag{5.6.14}
\end{equation}

\LayerTarget{tohoku:remarks:two-additional-spectral-sequences}
\noindent\textsc{Remarks.}\quad Let us record two further spectral sequences,
in two special cases, converging to $\operatorname{Ext}_{O,G}^n(X;A,B)$.

\LayerTarget{tohoku:construction:constant-equivariant-crossed-product}
Suppose that $O$ is the constant $G$-sheaf of rings defined by a ring $O$ on
which $G$ acts, and that $A$ is the constant $G$-$O$-sheaf $\mathbf M$
defined by an $O$-module $M$ on which $G$ acts compatibly with its operations
on $O$: $g^\circ(\lambda m)=g(\lambda)g(m)$. Introduce the ring $U$ generated
by $O$ and $\mathbf Z(G)$, subject to the commutation relations
$g\lambda g^{-1}=g(\lambda)$; thus $M$ is a $U$-module. Consider the functor
$h_{U,M}$ on the category of $U$-modules, defined by
$h_{U,M}(N)=\operatorname{Hom}_U(M,N)$. Then, for every $G$-$O$-module $B$,
$\operatorname{Hom}_{O,G}(M,B)=\operatorname{Hom}_U(M,\Gamma_X(B))$; that is,
one has the functorial identity

\LayerTarget{tohoku:equation:5.6.5-bis}
\begin{equation}
h_{O,G,M}=h_{U,M}\Gamma_X.
\tag{5.6.5 bis}
\end{equation}

\LayerTarget{tohoku:spectral-sequence:iv}
One easily checks further that $\Gamma_X$ takes injective objects of
$C^{O(G)}$ to injective $U$-modules. Theorem 2.4.1 therefore gives a
cohomological spectral functor IV on the category $C^{O(G)}$, converging to
the graded functor $\bigl(\operatorname{Ext}_{O,G}^n(X;M,B)\bigr)$, whose
initial term is

\LayerTarget{tohoku:equation:5.6.11-bis}
\begin{equation}
\mathrm {IV}_2^{p,q}(B)=\operatorname{Ext}_U^p\bigl(M,H^q(X,B)\bigr).
\tag{5.6.11 bis}
\end{equation}

\LayerTarget{tohoku:specialization:trivial-action-on-m}
If, in addition, $G$ acts trivially on $O$ and on $M$, then one also has
\[
\operatorname{Hom}_{O,G}(M,B)
=\operatorname{Hom}_U(M,\Gamma_X(B))
=\operatorname{Hom}_O(M,\Gamma_X^G(B)),
\]
whence the functorial formula

\LayerTarget{tohoku:equation:5.6.5-ter}
\begin{equation}
h_{O,G,M}=h_{O,M}\Gamma_X^G.
\tag{5.6.5 ter}
\end{equation}

\LayerTarget{tohoku:spectral-sequence:v}
One can further show that $\Gamma_X^G$ takes injective objects of $C^{O(G)}$
to injective $O$-modules (by showing that if $N$ is an injective $U$-module,
then $N^G$ is an injective $O$-module). Theorem 2.4.1 then gives a fifth
cohomological spectral functor, with the same abutment as the preceding ones,
whose initial term is

\LayerTarget{tohoku:equation:5.6.11-ter}
\begin{equation}
\mathrm V_2^{p,q}(B)=\operatorname{Ext}_O^p\bigl(M,H^q(X;G,B)\bigr).
\tag{5.6.11 ter}
\end{equation}

\LayerTarget{tohoku:comparison:spectral-sequences-iv-v}
If $G$ reduces to the identity, spectral sequences IV and V agree with the
spectral sequence of Theorem 4.3.1. If $X$ consists of a single point,
spectral sequence IV is trivial, whereas spectral sequence V is not trivial
in general and does not reduce to the preceding ones.\textsuperscript{12)}

\LayerTarget{tohoku:footnote:12}
\noindent\textsuperscript{12)} The initial term of this spectral sequence is
(with the notation of 11)) $\operatorname{Ext}_O^p\bigl(A,H^q(G,B)\bigr)$.

\TohokuContinuation{219}
